% Suprimir seccion y mover lo de Xu donde cuadre.
\subsection{Motivation: The Evaluation Gap in HGML}

% Revisar: Añadir referencia
Despite the proliferation of \gls{automl} tools, their systematic evaluation under \gls{hcai} criteria remains infrequent. Most comparative studies focus on predictive performance metrics, relegating to the background aspects such as usability for non-expert users, the transparency of automated decisions or the actual degree of control the system grants to the user \cite{hartikainen2022human}. Existing \gls{hgml} proposals \cite{gil2019towards,lindauer2024position} similarly lack evaluation instruments: they establish participation frameworks but provide no methodology for measuring the degree to which a tool actually realizes human-centered principles.

Applying \gls{hcai} criteria as an evaluative framework allows going beyond performance and answering questions such as: in which pipeline phases does the user have real intervention capacity?, which decisions remain opaque or beyond their reach?, how are responsibilities distributed between the human and the system? However, the \gls{hcai} frameworks available in the literature ---including Shneiderman's two-dimensional model \cite{shneiderman2022human}--- are formulated at a level of abstraction that does not translate directly into operational evaluation instruments for \gls{hgml} tools. They define desirable properties but do not specify \textit{how} to assess them in the context of a concrete \gls{ml} pipeline, \textit{which} phases to examine, or \textit{what} criteria determine whether a given decision genuinely belongs to the user or to the automated system.

This gap has been recently acknowledged at a broader level by Xu et al. \cite{xu2026hcai}, who propose a general \gls{hcai} methodological framework for \gls{ai} systems. However, their framework operates at a level of generality that does not address the specific characteristics of \gls{hgml} tools, leaving the need for a domain-specific evaluation methodology unresolved.

It is precisely because this gap between conceptual frameworks and operational evaluation methodology persists that the present work proposes a structured methodology that fills it: one grounded in the \gls{hgml} pipeline, operationalized through explicit responsibility assignment criteria, and applicable to any \gls{hgml} tool in a reproducible and comparable manner.